\section*{Capítulo \ref{ch:g11:infrared} --- A espectroscopia no infravermelho}

\begin{solution}{exo:g11:infrared:1}
$\lambda = 1/3300 = \qty{3.03e-4}{cm} = \qty{3.03}{\micro m}$;
$\lambda = 1/1050 = \qty{9.52e-4}{cm} = \qty{9.52}{\micro m}$. A banda em
\qty{3300}{cm^{-1}}, de \omterm{def:g11:infrared:wavenumber}{número de onda} maior, pertence à vibração mais
rápida.
\end{solution}

\begin{solution}{exo:g11:infrared:2}
$\qty{10}{\micro m} = \qty{1.0e-3}{cm}$, logo $\sigma = 1/\num{1.0e-3} =
\qty{1000}{cm^{-1}}$.
\end{solution}

\begin{solution}{exo:g11:infrared:3}
O \ce{C=O} de uma \omterm{def:g11:functional-groups:carbonyl}{cetona}; um \ce{O-H} de \omterm{def:g11:functional-groups:alcohol}{álcool} com \omterm{def:g11:polarity-and-cohesion:hydrogen-bond}{ligações de
hidrogênio}; ligações \ce{C-H} de uma cadeia de carbono.
\end{solution}

\begin{solution}{exo:g11:infrared:4}
A banda em \qty{1715}{cm^{-1}}: a ligação \ce{C=O}.
\end{solution}

\begin{solution}{exo:g11:infrared:5}
O instrumento mede a luz que atravessa a amostra; onde uma ligação
absorve, passa menos luz e a curva desce. $T =
\qty{100}{\percent}$ quer dizer que nenhuma luz desse \omterm{def:g11:infrared:wavenumber}{número de onda} é
absorvida.
\end{solution}

\begin{solution}{exo:g11:infrared:6}
O primeiro tem uma banda \ce{C=O} em \qty{1715}{cm^{-1}} e nenhuma banda
\ce{O-H}. É uma \omterm{def:g11:functional-groups:carbonyl}{cetona}, a butanona (também chamada butan-2-ona),
\ce{CH3-CO-CH2-CH3}. O segundo tem uma banda \ce{O-H} e nenhum \ce{C=O}: um
\omterm{def:g11:functional-groups:alcohol}{álcool}, por exemplo o butan-1-ol.
\end{solution}

\begin{solution}{exo:g11:infrared:7}
Os grupos \ce{O-H} do ácido têm \omterm{def:g11:polarity-and-cohesion:hydrogen-bond}{ligações de hidrogênio} muito fortes (as
\omterm{def:g7:atoms-and-molecules:molecule}{moléculas} se juntam aos pares, cada \ce{O-H} ligado ao \ce{C=O} da
outra), e cada \omterm{def:g7:atoms-and-molecules:molecule}{molécula} está ligada de um jeito um pouco diferente: a
banda fica ainda mais afrouxada e espalhada por uma faixa muito larga.
\end{solution}

\begin{solution}{exo:g11:infrared:8}
No líquido (A), a banda \ce{O-H} é larga, perto de
\qty{3350}{cm^{-1}}; no gás (B), onde as \omterm{def:g7:atoms-and-molecules:molecule}{moléculas} estão afastadas e não
formam \omterm{def:g11:polarity-and-cohesion:hydrogen-bond}{ligações de hidrogênio}, ela é uma banda fina perto de
\qty{3666}{cm^{-1}}.
\end{solution}

\begin{solution}{exo:g11:infrared:9}
Mal: perto de \qty{1730}{cm^{-1}} para o \omterm{def:g11:functional-groups:carbonyl}{aldeído} e de
\qty{1715}{cm^{-1}} para a \omterm{def:g11:functional-groups:carbonyl}{cetona}, valores próximos. A comparação do
espectro inteiro, \omterm{def:g11:infrared:band}{região da impressão digital} incluída, com espectros de
referência decidiria, ou outra técnica.
\end{solution}

\begin{solution}{exo:g11:infrared:10}
Um \omterm{def:g11:functional-groups:carboxylic-acid}{éster} (\ce{C=O} perto de \qty{1735}{cm^{-1}}, \ce{C-O} forte, nenhum
\ce{O-H}), por exemplo o etanoato de etila. O ácido etanoico é
\ce{C2H4O2}, e não \ce{C4H8O2}. O ácido butanoico tem a fórmula certa,
mas mostraria a banda \ce{O-H} muito larga dos ácidos: excluído.
\end{solution}

\begin{solution}{exo:g11:infrared:11}
O propan-1-ol é um \omterm{def:g11:functional-groups:alcohol}{álcool}: como o A. O ácido propanoico é um \omterm{def:g11:functional-groups:carboxylic-acid}{ácido
carboxílico}: como o D.
\end{solution}

\begin{solution}{exo:g11:infrared:12}
A banda \ce{O-H} larga perto de \qty{3350}{cm^{-1}} diminui enquanto uma
banda \ce{C=O} forte perto de \qty{1715}{cm^{-1}} cresce. A reação está
completa quando a banda \ce{O-H} desapareceu e a banda \ce{C=O} para de
crescer.
\end{solution}

\begin{solution}{exo:g11:infrared:13}
O ácido butanoico mostra uma banda \ce{O-H} muito larga de 2500 a
\qty{3300}{cm^{-1}}, o \omterm{def:g11:functional-groups:carboxylic-acid}{éster} nenhuma; o \ce{C=O} do ácido fica perto de
\qty{1710}{cm^{-1}}, mais largo, o do \omterm{def:g11:functional-groups:carboxylic-acid}{éster} perto de
\qty{1735}{cm^{-1}}, fino, com uma banda \ce{C-O} forte perto de
\qty{1240}{cm^{-1}}. O ácido tem um grupo \ce{O-H}, com \omterm{def:g11:polarity-and-cohesion:hydrogen-bond}{ligações de
hidrogênio}, que também afrouxa um pouco o seu \ce{C=O}.
\end{solution}

\begin{solution}{exo:g11:infrared:14}
Etanol que sobrou (o \omterm{def:g11:functional-groups:alcohol}{álcool} a partir do qual o \omterm{def:g11:functional-groups:carboxylic-acid}{éster} é produzido) ou
água (do \omterm{def:g4:air-a-mixture-of-gases:air}{ar} ou da lavagem): os dois têm grupos \ce{O-H} com \omterm{def:g11:polarity-and-cohesion:hydrogen-bond}{ligações de
hidrogênio}. A amostra pode ser comparada com um espectro de referência,
secada, destilada, ou ter a sua temperatura de ebulição medida.
\end{solution}

\begin{solution}{exo:g11:infrared:15}
Uma \omterm{def:g10:lewis-and-shape:multiple-bond}{ligação dupla} é uma mola mais dura que uma ligação simples entre os
mesmos \omterm{def:g7:atoms-and-molecules:atom}{átomos}: ela vibra mais depressa, num \omterm{def:g11:infrared:wavenumber}{número de onda} maior (1715
contra \qty{1050}{cm^{-1}}). As ligações \ce{O-H}, \ce{N-H} e \ce{C-H}
envolvem todas um \omterm{def:g7:atoms-and-molecules:atom}{átomo} de hidrogênio, o mais leve de todos: bolas leves
numa mola vibram depressa, daí \omterm{def:g11:infrared:wavenumber}{números de onda} perto de
2900--\qty{3600}{cm^{-1}}.
\end{solution}

\begin{solution}{pb:g11:infrared:1}
\textbf{1.} \ce{CH3-CH2-OH}; \ce{CH3-CO-CH3}; \ce{CH3-COOH}.

\textbf{2.} \ce{C2H6O}; \ce{C3H6O}; \ce{C2H4O2}.

\textbf{3.} Hidroxila; carbonila (\omterm{def:g11:functional-groups:carbonyl}{cetona}); carboxila.

\textbf{4.} Etanol: \ce{O-H} e \ce{C-H}. Propanona: \ce{C-H} e
\ce{C=O}. Ácido etanoico: \ce{O-H}, \ce{C-H} e \ce{C=O}.

\textbf{5.} Nenhuma banda \ce{C=O}, uma banda \ce{O-H} larga: o etanol.

\textbf{6.} Uma banda \ce{O-H} muito larga de 2500 a \qty{3300}{cm^{-1}}
e um \ce{C=O} forte perto de \qty{1710}{cm^{-1}}: o ácido etanoico.

\textbf{7.} A propanona: a banda \ce{C=O} muito forte em
\qty{1715}{cm^{-1}}, sem \ce{O-H}.

\textbf{8.} Provavelmente (vinagre, \omterm{def:g6:solutions-and-solubility:solute}{solvente}, \omterm{def:g11:functional-groups:alcohol}{álcool}), mas cheirar
líquidos desconhecidos é perigoso; o espectro é seguro e certo.

\textbf{9.} Os grupos \ce{O-H} do ácido têm \omterm{def:g11:polarity-and-cohesion:hydrogen-bond}{ligações de hidrogênio} mais
fortes (\omterm{def:g7:atoms-and-molecules:molecule}{moléculas} aos pares), por isso a banda fica mais baixa e muito
mais larga.

\textbf{10.} \ce{CH3-O-CH3}, \ce{C2H6O}; ele não tem \ce{O-H}, por isso
não tem banda \ce{O-H}.

\textbf{11.} A banda \ce{C=O}: perto de \qty{1730}{cm^{-1}} para o
propanal, perto de \qty{1715}{cm^{-1}} para a propanona.

\textbf{12.} A banda \ce{O-H}: um \omterm{def:g11:functional-groups:carboxylic-acid}{éster} não tem \ce{O-H}. O seu \ce{C=O}
ficaria perto de \qty{1735}{cm^{-1}}.

\textbf{13.} Os \omterm{def:g11:organic-skeletons:isomer}{isômeros} compartilham uma fórmula; os seus \omterm{def:g11:functional-groups:functional-group}{grupos
funcionais} diferem, e o espectro mostra os grupos.

\textbf{14.} \qty{1715}{cm^{-1}}.

\textbf{15.} $1715 \times 100 = \qty{1.715e5}{m^{-1}}$.

\textbf{16.} $\lambda = 1/\num{1.715e5} = \qty{5.83e-6}{m} =
\qty{5.83}{\micro m}$.

\textbf{17.} Não: a luz visível vai de cerca de 0.4 a
\qty{0.75}{\micro m}. É uma radiação infravermelha.

\textbf{18.} $\num{1.05e5}\ \si{m^{-1}}$, logo $\lambda = \qty{9.52e-6}{m}
= \qty{9.52}{\micro m}$.

\textbf{19.} A ligação \ce{C=O} (\omterm{def:g11:infrared:wavenumber}{número de onda} maior, comprimento de
onda mais curto).

\textbf{20.} Cerca de \qty{5.8}{\micro m}.
\end{solution}
